\begin{afterword}
\summaryhead

The post is written in the second person, much of it ironically. A TV pundit can explain any movement of bond yields after the fact, so the possible outcomes need not be treated as competing. A novice pundit, though, has 100 minutes to prepare excuses in advance for yields going up, down or staying the same, and must decide how to divide the time.

The pundit, who believes that probabilities are only ``little numbers'' in school word problems, finds that plausibility is no guide, since all three outcomes are plausible. But anticipation behaves differently from plausibility. It is limited: expecting one outcome more means expecting the others less. It seems to be related to how the time should be divided, and a parenthesis notes that with a logarithmic payoff the two can be matched. The post suggests treating anticipation as a conserved resource, like money, to be allocated. It mocks an alternative that uses ``pairs of numbers'' and ends by asking what an art of focusing uncertainty would be called.

\responsehead

The post makes one idea concrete, and makes it well: belief is a fixed quantity, so expecting one outcome more means expecting the others less. The pundit's 100 minutes turn the rule that probabilities add up to one into something a reader can feel. The contrast between limited anticipation and an unlimited ability to explain is also clear.

The post connects anticipation to time only in an aside. The post says, ironically, that the relation between them ``can't actually be quantified,'' and then supplies in a parenthesis the assumption that quantifies it: a payoff logarithmic in the time spent. Under that assumption the best plan gives each excuse minutes in proportion to its probability. Under others it does not. If the payoff grows in proportion to time, all 100 minutes go to the likeliest outcome. A reader who misses the parenthesis may take the match between time and belief as obvious. It holds under one assumption.

The post dismisses a rival theory with a pun. ``Pairs of numbers'' called ``DS'' is presumably Dempster--Shafer theory, whose lower and upper numbers are belief and plausibility. The reminder that follows, that the pundit has only 100 minutes, suggests that a fixed budget rules the theory out. It does not by itself, because that theory also divides a fixed total of 1, among sets of outcomes. The post may be right to prefer probabilities, but here it gives only the joke.

Finally, the post does not answer its closing question. The art it asks for exists: degrees of belief that obey the probability calculus and guide choices were worked out by Ramsey in 1926 and Savage in 1954, among others. The post names neither the theory nor its authors.

In short: a vivid case for treating belief as a fixed quantity to be divided; its one exact claim depends on an assumption stated only in a parenthesis, and it dismisses the rival theory with a pun.
\end{afterword}
